% Section: Human-in-the-Loop Governance
\section{Human-in-the-Loop Governance}
\label{sec:human-in-the-loop}

% TODO: Expand this section.
%
% Key topics:
% - Definition and role of human-in-the-loop (HITL) in recursive systems
% - When HITL checkpoints should be triggered
% - Checkpoint design principles
% - Minimizing cognitive load while maintaining meaningful oversight
% - Escalation protocols
% - Tooling and interface considerations

Human-in-the-loop (HITL) governance provides the essential alignment bridge between
autonomous agent behavior and human intent. In the Reflector framework, HITL governance
is not an ad-hoc override mechanism but a first-class architectural component with
well-defined trigger conditions, review interfaces, and approval protocols.

\subsection{Governance Checkpoints}

A \textit{governance checkpoint} is a defined point in the development workflow at
which autonomous progression is suspended pending human review and approval. Unlike
traditional code review, which is triggered by artifact production (e.g., a pull
request), Reflector governance checkpoints are triggered by \textit{system state events}:

\begin{itemize}
  \item Milestone completion
  \item Drift detection alerts
  \item Validation invariant violations
  \item Elapsed time thresholds
  \item Autonomous recursion depth limits
\end{itemize}

\subsection{Review Interface Design}

Effective HITL governance requires that human reviewers be presented with a clear,
concise summary of the system state at each checkpoint. The review interface must:

\begin{itemize}
  \item Summarize what the agent has done since the last checkpoint
  \item Highlight any deviations from the expected state
  \item Provide actionable approval and rejection options
  \item Surface the relevant validation results from the auditing system
  \item Minimize the time required for a meaningful review
\end{itemize}

\subsection{Approval Semantics}

An approval at a governance checkpoint carries a specific semantic meaning: the human
reviewer attests that the current system state is acceptable as a foundation for the
next autonomous phase. This approval is recorded in the governance audit trail and
serves as an immutable reference point for future drift analysis.
